Was vertrekken Gods wil
of een brutale gril

en volgde nu haar straf,
vroeg Emmy zich suf af.

De vraag bracht haar terug
naar de houten loopbrug

over een stuk moeras,
waarachter ze safe was.

Ook toen twijfelde ze
of haar fatalisme

door God was voorbestemd.
Of had ze hem ontstemd?

Haar vluchtweg werd plots steil,
afglijdend in onheil.

De klimmende helling
werd steeds meer een kwelling.

Langzaam verloor het wicht
haar wankel evenwicht.

Door een gebrek aan grip
raakte ze in een slip

de diepe afgrond in.
Beheersing had geen zin.

Ze raakte zichzelf kwijt
in de onzekerheid.

Een groep vogelvrijen,
weg van boerderijen,

nog verder van de stad,
kwam op Emmy’s zwerfpad.

Men vierde de cultus
van Petrus en Paulus

door witvis te eten
bij haar welkom heten.

Vriendelijke mensen
bedreigd met doodwensen,

daar men de benen nam.
Voor zichzelf eens opkwam.

Ze deelden hun boedel
en vormden een roedel.

Net als een groep herten
zag men in de verte

zo eerder een geldwolf.
Samen als een vloedgolf

konden ze die wel aan.
Er alleen voor te staan,

was een te groot gevaar.
Dat speelde niemand klaar.

De ontworpen keuring
neigde naar afwijzing

bleek uit zijn gelaatsscan.
“Weet je wel wie ik ben?”,

riep het aspirant lid
met verbeten gebit

tegen de beproever,
toetsend aan de oever

van een brede rivier.
Hij schreef op zijn papier

“Geestelijk instabiel:
de vorst is zo labiel

dat hij zichzelf niet kent.”
De arrogante vent,

die men al de bons gaf,
droop nogmaals onteerd af.

De beproever nam wraak
op iedere blaaskaak.

Hen die het lef hadden,
om hem bij zijn kladden

door terreur te pakken,
zette hij te kakken.

Emmy had als noviet,
die men zelden toeliet,

de laagste positie
in de groepshiërarchie.

Ze moest haar plaats weten.
Ze vernam hoe eten,

slapen, praten rechtaan
moesten worden gedaan.

Lege rituelen,
die Emmy moest delen.

Meespelen uit fatsoen. Waarom makkelijk doen

als het moeilijker kan?
Negeer je hersenpan

zoals een voetbalfan.
De vogelvrije clan

toonde haar allures
in haar procedures

over het toetreden.
De vijftien groepsleden

waren fysiek getest.
Dat beding bleek funest

te wezen voor het gros.
Men was al vaak de klos

als men moest gaan zwemmen
of zich moest afstemmen

op de groepsvluchtsnelheid.
Die raakte elk lijf kwijt

dat hen wilde vangen.
Dat kon sterk afhangen
van én de atmosfeer,
als ook de invloedssfeer

van de aardlichamen,
die eraan deelnamen.

Emmy miste de roep
van haar weglopersgroep,

dus reageerde traag.
Vogels bekeek ze graag

en dat werd haar verweer
tijdens haar terugkeer.

“Je houdt van vogelen?”
Het bewust goochelen

met dialectdwaling
gaf hen de voldoening

Emmy als een prooidier
voor erotisch plezier

gehaaid neer te zetten.
De misgreep rechtzetten

gaf hen nog meer kogels
“Nee, ik hou van vogels!

Ze maken mij zo blij,
want ze zijn zo gans vrij.”

Waarop men haar tartte:
“Je wilt dus van harte

met een vogel vrijen
dankzij vleierijen?”

Emmy dacht onvermoed:
“Ik versta hen niet goed”

en knikte dus maar ja.
Meteen volgde daarna

“Genoeg vogelvrijen
om mee te gaan vrijen!”

De mens als prooidier linkt
overlevingsinstinct

tegen roekeloosheid
aan angst en alertheid.

Waar Sigmund Freud op wees
was Emmy’s extra vrees.

Die vrijzwevende angst
tormenteerde het wrangst.

De beeldrijke woorden,
die haar hersens hoorden,

versterkten de dreiging
en haar vecht-vlucht-neiging.

Los van wetten en straf,
die Emmy’s vlucht haar gaf,

verdween haar slaafse angst.
Nochtans bleef ze het bangst,

hoewel als vogel vrij,
vanwege “timor Dei”,

haar heilige godsvrees,
die juist alhier verrees.

Haar geloof verliezen,
daar ze zelf kon kiezen,

tussen het goed of kwaad
in elke vrije daad.

Het in vrijheid leren,
ofwel het vrij leren,

over Emmy’s vrijheid,
spiritualiteit,

kerkelijke censuur,
en over de natuur

deed zijn klarende werk.
Emmy peinsde; “Ik merk

dat ik steeds verander”.
Al Anaximander

giste evolutie
als vrijheid in actie.

Anders dan hun beleid
was één lid ingewijd,

waarvan men vermoedde,
dat onder de hoede

van de roedelleider,
als zorgbegeleider,

hij was voorgetrokken.
Het was niet echt jokken,

maar erover zwijgen.
Het bleek de bloedeigen

broer te zijn van de baas.
Was er wellicht een maas

in de toetredingswet
en sprake van opzet?

De groep roemde verder
de alfa als “herder”,

omdat hij met zijn kracht
veiligheid had gebracht.

De herder kon zijn broer
zelfs niet als een tamboer

bij een leger lozen,
dus had hij gekozen

om hem op te nemen
ondanks zijn problemen.

Hij had een afwijking.
Waardoor één arm scheef hing.

Zijn tanden waren scheef.
Zijn pik lek als een zeef.

Zijn ogen keken schuin.
Geen haren op zijn kruin.

Mentaal minder begaafd
maakte hem onbeschaafd.

Er was niets aanvechtbaar
toerekeningsvatbaar

aan zijn inwendige.
Ook zijn uitwendige

kon je schrik bezorgen.
Emmy moest verzorgen.

Hem wassen, verplegen
en te eten geven.

Ze zorgde voor vermaak.
Dit alles was haar taak

als laagste roedellid.
Zijn vaste babysit.

De troep zag hem als beest.
Emmy als arm van geest.

Ze noemde hem “Einfach”,
omdat ze zijn aard zag.

Was ze misschien verblind
door dit duivelse kind?
